\section{BORNOLOGICAL SPACES}

In this paragraph, E denotes a locally convex space, and B its canonical bornology (III, p.~3, def. 5).

Lemma 1. --- Let G be a semi-normed space, p its semi-norm, and u a linear mapping from G into E. The following conditions are equivalent :

\begin{enumerate}
\def\labelenumi{(\roman{enumi})}
\item
  \(u\) is continuous;
\item
  the image of the unit ball of \(G\) under \(u\) is bounded in \(\mathbf{E}\) ;
\item
  for every sequence \((x_{n})\) of points of \(\mathbf{G}\) tending to 0, the sequence \((u(x_{n}))\) is bounded in E.
\end{enumerate}

It is immediate that (i) implies (ii) (III, p.~4, cor. 1) and that (ii) implies (iii). Let V be a neighbourhood of 0 in E; if \(u^{-1}(\mathrm{V})\) is not a neighbourhood of 0 in G, then there exists a sequence \((y_n)\) of points of \(\mathrm{G} - u^{-1}(\mathrm{V})\) such that \(p(y_n) \leqslant \frac{1}{n^2}\) . Hence the sequence \(x_n = ny_n\) tends to 0 in G and \(u(x_n) \notin n\mathrm{V}\) , which implies that the sequence \((u(x_n))\) is not bounded. Therefore (iii) implies (i).

PROPOSITION 1. -- The following conditions are equivalent :

\begin{enumerate}
\def\labelenumi{(\roman{enumi})}
\tightlist
\item
  Every semi-norm on \(\mathbf{E}\) which is bounded on bounded subsets of \(\mathbf{E}\) is continuous.
\end{enumerate}

(i') Every convex balanced subset of E which absorbs the bounded subsets of E (I, p.~7, def. 4) is a neighbourhood of 0 in E.

\begin{enumerate}
\def\labelenumi{(\roman{enumi})}
\setcounter{enumi}{1}
\tightlist
\item
  E is the inductive limit of the semi-normed spaces \(E_{A}\) , where A ranges over the directed increasing set of closed, convex, balanced and bounded subsets of E.
\end{enumerate}

(ii') There exists a family \((\mathrm{E}_{i})_{i\in\mathrm{I}}\) of semi-normed spaces, and for every \(i\in I\) , a linear mapping \(u_{i}:E_{i}\to E\) such that the topology of E is the finest locally convex topology for which the \(u_{i}\) are continuous.

\begin{enumerate}
\def\labelenumi{(\roman{enumi})}
\setcounter{enumi}{2}
\tightlist
\item
  For an arbitrary locally convex space F, a linear mapping \(u: E \to F\) is continuous if and only if for every sequence \((x_{n})\) of points in E tending to 0, the sequence \((u(x_{n}))\) is bounded in F.
\end{enumerate}

(iii') For an arbitrary semi-normed space \(\mathbf{F}\) , a linear mapping \(u: \mathbf{E} \to \mathbf{F}\) is continuous if and only if \(u(\mathbf{X})\) is bounded in \(\mathbf{F}\) for every bounded set \(\mathbf{X}\) in \(\mathbf{E}\) .

It is immediate that (i) and (i') are equivalent in view of the correspondence between semi-norms and convex, balanced, absorbent subsets (II, p.~20). If \(p\) is a semi-norm on E, which is continuous on each \(\mathrm{E}_{\mathrm{A}}\) , then \(p\) is bounded on bounded subsets of E; hence (i) implies (ii) (II, p.~27, prop. 5). It is clear that (ii) implies (ii').

Now let \((\mathrm{E}_{i}, u_{i})_{i\in I}\) be as in (ii') and let u be a linear mapping from E into a locally convex space F, such that \((u(x_{n}))\) is bounded in F for every sequence \((x_{n})\) of points of E tending to 0. It follows from lemma 1 of III, p.~11 that the linear mapping \(u\circ u_{i}:E_{i}\to F\) is continuous for all \(i\in I\) . Hence, if the topology of E is the finest locally convex topology for which the \(u_{i}\) are continuous, then u is continuous (II, p.~27, prop. 5). This proves that (ii') implies (iii).

It is immediate that (iii) implies (iii') (III, p.~3, cor.) Finally, if \(p\) is a semi-norm on E, which is bounded on bounded subsets of E, the condition (iii') asserts that the identity map is continuous from E into the semi-normed space (E, \(p\) ); in other words, \(p\) is continuous. This proves that (iii') implies (i).

DEFINITION 1. --- A locally convex space is said to be bornological if it satisfies the equivalent conditions of prop. 1.

Examples. --- 1) Every semi-normed space is bornological.

\begin{enumerate}
\def\labelenumi{\arabic{enumi})}
\setcounter{enumi}{1}
\item
  In particular, every finite dimensional locally convex space is bornological.
\item
  On account of the transitivity of final locally convex topologies (II, p.~28, cor. 2), we deduce at once from condition (ii') that if \((\mathrm{E}_i)_{i\in I}\) is a family of locally convex bornological spaces and if E is assigned the finest locally convex topology for which the linear mappings \(u_{i}:\mathrm{E}_{i}\to \mathrm{E}\) (for \(i\in \mathrm{I}\) ) are continuous, then E is bornological. In particular, an inductive limit, a direct sum, a quotient space of bornological spaces are bornological spaces.
\end{enumerate}

On the other hand, a closed subspace of a bornological space is not necessarily bornological (IV, p.~63, exerc. 8).

COROLLARY. --- Every Hausdorff and semi-complete bornological space is an inductive limit of Banach spaces.

In fact, the spaces \(E_{A}\) , where A is closed and bounded are Banach spaces (III, p.~8, corollary).

PROPOSITION 2. --- A locally convex metrizable space is bornological.

Suppose E is metrizable, and p a semi-norm on E which is bounded on bounded subsets of E, but which is not continuous. Let A be the set of all \(x \in E\) such that \(p(x) < 1\) . Let \((\mathrm{V}_{n})_{n \geqslant 1}\) be a decreasing sequence forming a fundamental system of neighbourhoods of 0 in E. Since p is not continuous, A is not a neighbourhood of 0; hence for every n \textgreater{} 0, we have \(A \not\subset n^{-1}V_{n}\) and there exists a point \(x_{n}\) in \(V_{n}\) , such that \(n^{-1}x_{n} \notin A\) , that is, \(p(x_{n}) \geqslant n\) . The sequence \((x_{n})\) tends to 0, hence is bounded (III, p.~3, corollary); this contradicts the hypothesis on p.

COROLLARY. --- Every Fréchet space (II, p.~24) is the inductive limit of Banach spaces.

